%% ut-thesis.tex -- document template for graduate theses at UofT
%%
%% Copyright (c) 1998-2013 Francois Pitt <fpitt@cs.utoronto.ca>
%% last updated at 16:20 (EDT) on Wed 25 Sep 2013
%%
%% This work may be distributed and/or modified under the conditions of
%% the LaTeX Project Public License, either version 1.3c of this license
%% or (at your option) any later version.
%% The latest version of this license is in
%%     http://www.latex-project.org/lppl.txt
%% and version 1.3c or later is part of all distributions of LaTeX
%% version 2005/12/01 or later.
%%
%% This work has the LPPL maintenance status "maintained".
%%
%% The Current Maintainer of this work is
%% Francois Pitt <fpitt@cs.utoronto.ca>.
%%
%% This work consists of the files listed in the accompanying README.

%% SUMMARY OF FEATURES:
%%
%% All environments, commands, and options provided by the `ut-thesis'
%% class will be described below, at the point where they should appear
%% in the document.  See the file `ut-thesis.cls' for more details.
%%
%% To explicitly set the pagestyle of any blank page inserted with
%% \cleardoublepage, use one of \clearemptydoublepage,
%% \clearplaindoublepage, \clearthesisdoublepage, or
%% \clearstandarddoublepage (to use the style currently in effect).
%%
%% For single-spaced quotes or quotations, use the `longquote' and
%% `longquotation' environments.


%%%%%%%%%%%%         PREAMBLE         %%%%%%%%%%%%

%%  - Default settings format a final copy (single-sided, normal
%%    margins, one-and-a-half-spaced with single-spaced notes).
%%  - For a rough copy (double-sided, normal margins, double-spaced,
%%    with the word "DRAFT" printed at each corner of every page), use
%%    the `draft' option.
%%  - The default global line spacing can be changed with one of the
%%    options `singlespaced', `onehalfspaced', or `doublespaced'.
%%  - Footnotes and marginal notes are all single-spaced by default, but
%%    can be made to have the same spacing as the rest of the document
%%    by using the option `standardspacednotes'.
%%  - The size of the margins can be changed with one of the options:
%%     . `narrowmargins' (1 1/4" left, 3/4" others),
%%     . `normalmargins' (1 1/4" left, 1" others),
%%     . `widemargins' (1 1/4" all),
%%     . `extrawidemargins' (1 1/2" all).
%%  - The pagestyle of "cleared" pages (empty pages inserted in
%%    two-sided documents to put the next page on the right-hand side)
%%    can be set with one of the options `cleardoublepagestyleempty',
%%    `cleardoublepagestyleplain', or `cleardoublepagestylestandard'.
%%  - Any other standard option for the `report' document class can be
%%    used to override the default or draft settings (such as `10pt',
%%    `11pt', `12pt'), and standard LaTeX packages can be used to
%%    further customize the layout and/or formatting of the document.

%% *** Add any desired options. ***


\documentclass[11pt]{ut-thesis}

\usepackage{appendix}
\usepackage{graphicx}
\usepackage{amsmath}
\usepackage{amssymb}
\usepackage{amsthm}
\usepackage{mathrsfs}
\usepackage[byname]{smartref}
\usepackage{hyperref} %comment out for hardcopy
%\usepackage{txfonts}
\usepackage{tocloft}
\usepackage[numbers]{natbib}
\usepackage[symbols,nogroupskip,nonumberlist]{glossaries-extra}
\usepackage{setspace}
\usepackage{blkarray}

\newcommand{\matindex}[1]{\mbox{\scriptsize#1}}

\makeglossaries

\glsxtrnewsymbol[description={Individual index}]{i}{\ensuremath{i}}
\glsxtrnewsymbol[description={Family (Cluster) index}]{j}{\ensuremath{j}}
\glsxtrnewsymbol[description={Proband index}]{p}{\ensuremath{p}}
\glsxtrnewsymbol[description={Number of events in family $j$}]{dj}{$d_j$}
\glsxtrnewsymbol[description={Some time}]{t}{\ensuremath{t}}
\glsxtrnewsymbol[description={Some Time for the proband}]{a}{\ensuremath{a}}
\glsxtrnewsymbol[description={Event Time}]{T}{\ensuremath{T}}
\glsxtrnewsymbol[description={Event indicator for individual $i$ in family $j$}]{delta}{$\delta_{ij}$}
\glsxtrnewsymbol[description={The observed survival data ($t,\delta$)}]{w}{\ensuremath{w}}
\glsxtrnewsymbol[description={Number of individuals}]{n}{\ensuremath{n}}
\glsxtrnewsymbol[description={Number of Families (Clusters)}]{J}{\ensuremath{J}}
\glsxtrnewsymbol[description={Index of the sampled completed dataset in the MCEM}]{m}{\ensuremath{m}}
\glsxtrnewsymbol[description={Number of the sampled completed dataset in the MCEM}]{M}{\ensuremath{M}}
\glsxtrnewsymbol[description={Frailty term}]{z}{\ensuremath{z}}
\glsxtrnewsymbol[description={$q$-th element of Gauss Hermite Quadrature}]{q}{\ensuremath{q}}
\glsxtrnewsymbol[description={$q$-th weight of Gauss Hermite Quadrature}]{omega}{\ensuremath{\omega}}
\glsxtrnewsymbol[description={$q$-th node of Gauss Hermite Quadrature}]{yq}{\ensuremath{y_q}}
\glsxtrnewsymbol[description={Total number of quadratures}]{Nq}{\ensuremath{N_q}}
\glsxtrnewsymbol[description={Hazard fucntion}]{h}{\ensuremath{h(\cdot)}}
\glsxtrnewsymbol[description={Baseline hazard function}]{h_0}{\ensuremath{h_0(\cdot)}}
\glsxtrnewsymbol[description={Cumulative hazard fucntion}]{H}{\ensuremath{H(\cdot)}}
\glsxtrnewsymbol[description={Survival fucntion}]{S}{\ensuremath{S(\cdot)}}
\glsxtrnewsymbol[description={Ascertainment of family $j$ into the study}]{A}{$A_j(\cdot)$}
\glsxtrnewsymbol[description={Likelihood function}]{L}{$L(\cdot)$}
\glsxtrnewsymbol[description={Log-likelihood function}]{ell}{$\ell(\cdot)$}
\glsxtrnewsymbol[description={Laplace transform}]{lap}{$\mathscr{L}(\cdot)$}
\glsxtrnewsymbol[description={Covariates}]{X}{\ensuremath{\mathbf{x}}}
\glsxtrnewsymbol[description={Model coefficients vector}]{beta}{$\boldsymbol{\beta}$}
\glsxtrnewsymbol[description={Parameter vector}]{theta}{$\boldsymbol{\theta}$}
\glsxtrnewsymbol[description={The combination of $(\boldsymbol{\beta}, \lambda, \alpha)$}]{Lambda}{$\Lambda$}
\glsxtrnewsymbol[description={Weibull shape parameter}]{lambda}{$\lambda$}
\glsxtrnewsymbol[description={Weibull scale parameter}]{alpha}{$\alpha$}
\glsxtrnewsymbol[description={General form of the parameter in an undefined frailty distribution}]{upsilon}{$\upsilon$}
\glsxtrnewsymbol[description={Gamma shape and rate parameters}]{k}{$k$}
\glsxtrnewsymbol[description={Log-Normal variance parameter}]{sigma}{$\sigma^2$}
\glsxtrnewsymbol[description={Missing data distribution parameters}]{psi}{$\psi$}



%% *** Add \usepackage declarations here. ***
%% The standard packages `geometry' and `setspace' are already loaded by
%% `ut-thesis' -- see their documentation for details of the features
%% they provide.  In particular, you may use the \geometry command here
%% to adjust the margins if none of the ut-thesis options are suitable
%% (see the `geometry' package for details).  You may also use the
%% \setstretch command to set the line spacing to a value other than
%% single, one-and-a-half, or double spaced (see the `setspace' package
%% for details).


%%%%%%%%%%%%%%%%%%%%%%%%%%%%%%%%%%%%%%%%%%%%%%%%%%%%%%%%%%%%%%%%%%%%%%%%
%%                                                                    %%
%%                   ***   I M P O R T A N T   ***                    %%
%%                                                                    %%
%%  Fill in the following fields with the required information:       %%
%%   - \degree{...}       name of the degree obtained                 %%
%%   - \department{...}   name of the graduate department             %%
%%   - \gradyear{...}     year of graduation                          %%
%%   - \author{...}       name of the author                          %%
%%   - \title{...}        title of the thesis                         %%
%%%%%%%%%%%%%%%%%%%%%%%%%%%%%%%%%%%%%%%%%%%%%%%%%%%%%%%%%%%%%%%%%%%%%%%%

%% *** Change this example to appropriate values. ***
\degree{Master of Science}
\department{Epidemiology and Biostatistics\\
Schulich School of Medicine \& Dentistry}
\gradyear{2024}
\author{Jiaqi Bi}
\title{Frailty Model Incorporating Ascertainment Correction with Missing Data in Family-Based Study}

%% *** NOTE ***
%% Put here all other formatting commands that belong in the preamble.
%% In particular, you should put all of your \newcommand's,
%% \newenvironment's, \newtheorem's, etc. (in other words, all the
%% global definitions that you will need throughout your thesis) in a
%% separate file and use "\input{filename}" to input it here.


%% *** Adjust the following settings as desired. ***

%% List only down to subsections in the table of contents;
%% 0=chapter, 1=section, 2=subsection, 3=subsubsection, etc.
\setcounter{tocdepth}{2}

%% Make each page fill up the entire page.
\flushbottom


%%%%%%%%%%%%      MAIN  DOCUMENT      %%%%%%%%%%%%

\begin{document}


%% This sets the page style and numbering for preliminary sections.
\begin{preliminary}

%% This generates the title page from the information given above.
\maketitle

%% There should be NOTHING between the title page and abstract.
%% However, if your document is two-sided and you want the abstract
%% _not_ to appear on the back of the title page, then uncomment the
%% following line.
\cleardoublepage

%% This generates the abstract page, with the line spacing adjusted
%% according to SGS guidelines.
\begin{abstract}
    \addcontentsline{toc}{chapter}{Abstract}
    This is an abstract section...
%% *** Put your Abstract here. ***
%% (At most 150 words for M.Sc. or 350 words for Ph.D.)
\end{abstract}

%% Anything placed between the abstract and table of contents will
%% appear on a separate page since the abstract ends with \newpage and
%% the table of contents starts with \clearpage.  Use \cleardoublepage
%% for anything that you want to appear on a right-hand page.

%% This generates a "dedication" section, if needed -- just a paragraph
%% formatted flush right (uncomment to have it appear in the document).
\begin{dedication}
    \addcontentsline{toc}{chapter}{Dedication}
    \textit{To my mother, Kai Hua, for all the upbringrings. }

    \textit{To my father, Guangmin Bi, for teaching me not to give up. }
    
    \textit{To my supervisors, for their unstopping guidance and invaluable lessons. }

    \textit{To my friends, for their peer pressures and supports, and their companionships.}
%% *** Put your Dedication here. ***
\end{dedication}

%% The `dedication' and `acknowledgements' sections do not create new
%% pages so if you want the two sections to appear on separate pages,
%% uncomment the following line.
\newpage  % separate pages for dedication and acknowledgements

%% Alternatively, if you leave both on the same page, it is probably a
%% good idea to add a bit of extra vertical space in between the two --
%% for example, as follows (adjust as desired).
%\vspace{.5in}  % vertical space between dedication and acknowledgements

%% This generates an "acknowledgements" section, if needed
%% (uncomment to have it appear in the document).
\begin{acknowledgements}
    \addcontentsline{toc}{chapter}{Acknowledgements}
%% *** Put your Acknowledgements here. ***
\end{acknowledgements}

\newpage

%% This generates the Table of Contents (on a separate page).
\tableofcontents
\printunsrtglossary[type=symbols,style=long,title={List of Symbols}]

\newpage
%% This generates the List of Tables (on a separate page), if needed
%% (uncomment to have it appear in the document).
\listoftables

\newpage
%% This generates the List of Figures (on a separate page), if needed
%% (uncomment to have it appear in the document).
\listoffigures

%% You can add commands here to generate any other material that belongs
%% in the head matter (for example, List of Plates, Index of Symbols, or
%% List of Appendices).

%% End of the preliminary sections: reset page style and numbering.
\end{preliminary}


%%%%%%%%%%%%%%%%%%%%%%%%%%%%%%%%%%%%%%%%%%%%%%%%%%%%%%%%%%%%%%%%%%%%%%%%
%%  Put your Chapters here; the easiest way to do this is to keep     %%
%%  each chapter in a separate file and `\include' all the files.     %%
%%  Each chapter file should start with "\chapter{ChapterName}".      %%
%%  Note that using `\include' instead of `\input' will make each     %%
%%  chapter start on a new page, and allow you to format only parts   %%
%%  of your thesis at a time by using `\includeonly'.                 %%
%%%%%%%%%%%%%%%%%%%%%%%%%%%%%%%%%%%%%%%%%%%%%%%%%%%%%%%%%%%%%%%%%%%%%%%%
\setstretch{1.4}
\chapter{Time series: Long memory}
Here is a picture of a long memory time series. 
\begin{figure}[ht]
\begin{center}
%\includegraphics[height = 9cm, width = 9cm]{pic1.jpeg}
\includegraphics[height = 9cm, width = 9cm]{logo.png}
\caption{A long memory time series\label{ts1}}
\end{center}
\end{figure}

Here's a table.
\begin{table}[ht]
\begin{center}
\begin{tabular}[ht]{|c|lr|c|} 
%c stands for centre, l for left, r for right; the | puts lines in between, and the hline puts a horizontal line in
\hline
$n$ & $\alpha$ &$n\alpha$ & $\beta$\\
\hline
1 & 0.2 & 0.2 & 5\\
\hline
2 & 0.3 & 0.6 & 4\\
\hline
3 & 0.7 & 2.1 & 3\\
\hline
\end{tabular}
\caption{A random table \label{tab1}}
\end{center}
\end{table}

\begin{eqnarray}
y &=& mx + b \label{eq1}\\
&=& ax+ c
\label{eq2}
\end{eqnarray}

This is an un-numbered equation, along with a numbered one. 
\begin{eqnarray}
u &=& px \nonumber\\
p &=& P(X=x) \label{eqn3}
\end{eqnarray}

Look at Table \ref{tab1} and Figure \ref{ts1} and equations \ref{eq1},  \ref{eq2}, and \ref{eqn3}.

Let's do some matrix algebra now.

\begin{equation}
det\left(\left|\begin{array}{ccc} 2 & 3 & 5\\
4 & 4 & 6\\
9 & 8 & 1
\end{array}\right|\right) = 42
\end{equation}

In the equation and eqnarray environments, you don't need to have the dollar sign to enter math mode.

\begin{eqnarray}
\alpha = \beta_1 \Gamma^{-1}
\end{eqnarray}

This is citing a reference ~\cite{mygood11111}.  This is citing another ~\cite{mrx05}.  Nobody said something ~\cite{Nobody06}.

\chapter{Theorems}
\section{Basic Theorems}
\begin{theorem}
 $e^{i\pi} = -1$ \label{eipi}
\end{theorem}
\chapter{Methods}
\section{Introduction}
Missing data is a common and unavoidable topic for many studies that involved with data. 
Multiple imputation is a well-established method for dealing with the missing data in various statistical analyses as well. 
Until today, to our best knowledge, the multiple imputation in the context of frailty models is rarely discussed. 
The oversight is even more pronounced within the field of genetic epidemiology, such that the frailty model is added a layer of ascertainment correction term to account for selection bias that could be introduced in the data sampling stage.
Although it is important that many research may encounter the missing data, the current existing literature lacks comprehensive discussions and methodological developments that integrate multiple imputation with the uniqueness posed by frailty models within the genetic epidemiology area. 
Addressing the gap and proposing the method will enhance the validity of findings for this field. 

When imputing continuous variables, a common approach is to assume a conditionally normal distribution, which can be estimated from a linear regression.
The original multiple imputation framework, introduced by \citet{rubin1987multiple}, has been widely adopted by researchers across various scientific disciplines using different models.
There are many softwares that implement the multiple imputations, such as Blimp software~\cite{keller2021blimp}, \texttt{MICE} package in R~\cite{royston2011multiple}, and \texttt{jomo} package in R~\cite{quartagno2019jomo}. 
Apparently, due to space limitations, this thesis cannot discuss all the well-designed software available.
The development of Multiple Imputation within certain specific models remains incomplete, and many details have not yet been comprehensively addressed in the current published literature.
In genetic epidemiology, studies often involve family clusters, yet current multiple imputation methods do not account for genetic correlations.
Additionally, the application for multiple imputation combined with the frailty model with ascertainment correction has not been thoroughly explored.
Therefore, this research aims to develop a computationally efficient multiple imputation method that accounts for kinship correlations and applies it to frailty models with ascertainment correction.

In this chapter, a comprehensive guideline and adjusted multiple imputation formulas are provided. 
We explicitly demonstrate how the kinship matrix functions in the imputation step and how ascertainment correction is applied in the analysis step of this research.
Moreover, the specialized imputation model is introduced, taking into account genetic variance and residual variance (also known as environmental variance).
Then the variance estimation using Rubin's Rule is presented, along with the confidence interval based on the completed data following the proposed multiple imputation method.

\section{Kinship Matrix}
%\raggedbottom
The kinship matrix, also known as the relatedness matrix, is a fundamental concept in statistical genetics, particularly in the context of quantitative genetics and genetic epidemiology. 
It quantifies the genetic relatedness between individuals based on pedigree information~\cite{lynch1998genetics}. 
This matrix is essential for estimating heritability and for controlling for familial relatedness in genetic association studies. 
The elements of the kinship matrix represent the probability that a randomly chosen allele from one individual is identical by descent (IBD) to a randomly chosen allele from another individual.
Technically, the kinship matrix can be generated either for each family or for all individuals in a study, once the relationships between each individual in a family or study are known~\cite{sinnwell2014kinship2}.
For example, the kinship matrix $K$ for a family of four individuals may look like this:
\begin{equation*}
    K = \begin{bmatrix}
    1    & 0.5  & 0.25 & 0 \\
    0.5  & 1    & 0.5  & 0 \\
    0.25 & 0.5 & 1    & 0 \\
    0    & 0    & 0    & 1 \\
    \end{bmatrix}
\end{equation*}
Here, the rows and columns represent the same individuals.
The diagonal of the kinship matrix has a value of 1, indicating the correlation of an individual with themselves.
Specifically, 
\begin{enumerate}
    \item In $K_{11}$, the individual index $1$ and himself is fully related to himself. 
    \item In $K_{12}=0.5$, the first and second individuals are half-related (e.g., Siblings). 
    \item In $K_{13}=0.25$, the first and third individuals are a quarter-related (e.g., half-siblings or grandparent-grandchild). 
    \item In $K_{14}=0$, the first and the fourth individuals are not related (e.g., spousal relationship, or lawful adoption). 
\end{enumerate}

One advantage of using the kinship matrix is that whether the data is partitioned by family or considered globally for all individuals from different families, the results are mostly equivalent.
Below shows another example of the kinship matrix that how it looks like on a global setting for three families.
\begin{enumerate}
    \item For family 1, A is the parent, B is the child, C is the sibling of A, and D is the cousin of B. 
    \item For family 2, E is the parent, F is the child, G is the grandparent of F, H is the half-sibling of F
    \item For family 3, I is the parent, J is the child, K is the aunt or uncle of J.
\end{enumerate}
Among these three families, the matrix can be partitioned into three segments to represent the kinship correlations.
In practice, this is often unnecessary since unrelated individuals will have a correlation value of 0 in the matrix.
Here is an example matrix combining all families:

\begin{equation*}
    \scriptsize
    K = \begin{blockarray}{cccccccccccc}
        & \matindex{A} & \matindex{B} & \matindex{C} & \matindex{D} & \matindex{E} & \matindex{F} & \matindex{G} & \matindex{H} & \matindex{I} & \matindex{J} & \matindex{K} \\
        \begin{block}{c(ccccccccccc)}
            \matindex{A} & 1 & 0.5 & 0.5 & 0.125 & 0 & 0 & 0 & 0 & 0 & 0 & 0  \\
            \matindex{B} & 0.5 & 1 & 0.25 & 0.0625 & 0 & 0 & 0 & 0 & 0 & 0 & 0  \\
            \matindex{C} & 0.5 & 0.25 & 1 & 0.0625 & 0 & 0 & 0 & 0 & 0 & 0 & 0  \\
            \matindex{D} & 0.125 & 0.0625 & 0.0625 & 1 & 0 & 0 & 0 & 0 & 0 & 0 & 0  \\
            \matindex{E} & 0 & 0 & 0 & 0 & 1 & 0.5 & 0.25 & 0.125 & 0 & 0 & 0  \\
            \matindex{F} & 0 & 0 & 0 & 0 & 0.5 & 1 & 0.5 & 0.25 & 0 & 0 & 0  \\
            \matindex{G} & 0 & 0 & 0 & 0 & 0.25 & 0.5 & 1 & 0.125 & 0 & 0 & 0  \\
            \matindex{H} & 0 & 0 & 0 & 0 & 0.125 & 0.25 & 0.125 & 1 & 0 & 0 & 0  \\
            \matindex{I} & 0 & 0 & 0 & 0 & 0 & 0 & 0 & 0 & 1 & 0.5 & 0.25  \\
            \matindex{J} & 0 & 0 & 0 & 0 & 0 & 0 & 0 & 0 & 0.5 & 1 & 0.25  \\
            \matindex{K} & 0 & 0 & 0 & 0 & 0 & 0 & 0 & 0 & 0.25 & 0.25 & 1  \\
        \end{block}
    \end{blockarray}
\end{equation*}

The kinship matrix is meant to be symmetric, yielding real eigenvalues and orthogonal eigenvectors, which is beneficial in principal component analysis (PCA) often conducted in genome-wide association studies (GWAS) to reduce dimensionality~\cite{price2006principal}.
Also, this symmetric property ensures the genetic relatedness between any two individuals is consistent regardless of their order in a dataset, which further ensures an accurate interpretation in the the statistical inference and valid and unbiased statistical tests. 
In statistical analyses, the kinship matrix is not typically estimated; rather, it is derived from a precise understanding of familial relationships. 
This direct acquisition ensures an accurate representation of genetic relatedness, which is essential for the integrity of subsequent genetic analyses.

In some variables within a familial study, the structure could be highly dependent on the kinship matrix. 
For example, the polygenic risk score (PRS) may be used as a covariate when modeling certain disease risks.
Particularly in genetic studies with missing data, neglecting the kinship correlation structure can lead to biased parameter inference.
Current imputation models face the challenge of not adequately handling kinship correlations.
Thus, the following sections demonstrate the comparisons when conducting proper multiple imputation with and without ignoring pedigree information.

\section{Multiple Imputation without Kinship Matrix}
%\raggedbottom
Multiple imputation is an advantageous method for addressing missing data under the missing at random (MAR) mechanism, as it preserves the inherent variability of the data and produces valid statistical inferences. 
The imputation step appropriately accounts for uncertainties that may be caused by the missing data. 
This method has been proven to be robust for common statistical analyses, especially in contexts such as clinical trials or psychological observational studies where genetic elements are not present. 
Existing multiple imputation methods for continuous variables often assume a conditionally normal distribution. 
Suppose $y$ is the variable containing missing values, and $\mathbf{x}$ are other fully observed variables.
We assume there are total $p$ parameters being estimated in this linear regression. 
The conditionally normal model (linear regression) can be written as 
\begin{equation} 
    y|\mathbf{x},\boldsymbol{\beta}\sim N(\boldsymbol{\beta}\mathbf{x}, \sigma^2)
\end{equation}
where $\boldsymbol{\beta}$ here represents the true parameter value that brings a linearly associated relationship between observed $\mathbf{x}$ and $y$.
The $\sigma$ is the standard error introduced in this imputation model. 
In the linear regression setting, 
\begin{equation} 
    y=\mathbf{x}\boldsymbol{\beta}+\epsilon,~ \epsilon\sim N(0, \sigma^2)
\end{equation}
This simply corresponds to the likelihood 
\begin{equation}\label{eq:bayeslinearnormal}
    f(y|\mathbf{x}, \boldsymbol{\beta}, \sigma^2)\propto (\sigma^2)^{-n/2}\exp \Big (-\frac{1}{2\sigma^2} (y-\mathbf{x}\boldsymbol{\beta})^{\top}(y-\mathbf{x}\boldsymbol{\beta})\Big )
\end{equation}
because of the conditional normality of $y$. 
From the likelihood~\ref{eq:bayeslinearnormal}, we can solve that 
\begin{equation} 
    \hat{\boldsymbol{\beta}}=(\mathbf{x}^{\top}\mathbf{x})^{-1}\mathbf{x}^{\top}y
\end{equation}
%In the Bayesian framework, we need to find the prior, which is the joint distribution $f(\sigma^2,\boldsymbol{\beta})$. 
From~\ref{eq:bayeslinearnormal}, note that $(y-\mathbf{x}\boldsymbol{\beta})^{\top}(y-\mathbf{x}\boldsymbol{\beta})$ can be written as 
\begin{align} 
    (y-\mathbf{x}\boldsymbol{\beta})^{\top}(y-\mathbf{x}\boldsymbol{\beta})&=\big ((y-\mathbf{x}\boldsymbol{\beta})+(\mathbf{x}\hat{\boldsymbol{\beta}}-\mathbf{x}\boldsymbol{\beta})\big )^{\top}\big ((y-\mathbf{x}\boldsymbol{\beta})+(\mathbf{x}\hat{\boldsymbol{\beta}}-\mathbf{x}\boldsymbol{\beta})\big )\\
    &=(y-\mathbf{x}\boldsymbol{\beta})^{\top}(y-\mathbf{x}\boldsymbol{\beta}) + (\boldsymbol{\beta}-\hat{\boldsymbol{\beta}})^{\top}(\mathbf{x}^{\top}\mathbf{x})(\boldsymbol{\beta}-\hat{\boldsymbol{\beta}})+2(\mathbf{x}\hat{\boldsymbol{\beta}}-\mathbf{x}\boldsymbol{\beta})\\
    &=(y-\mathbf{x}\boldsymbol{\beta})^{\top}(y-\mathbf{x}\boldsymbol{\beta}) + (\boldsymbol{\beta}-\hat{\boldsymbol{\beta}})^{\top}(\mathbf{x}^{\top}\mathbf{x})(\boldsymbol{\beta}-\hat{\boldsymbol{\beta}})
\end{align}
So Equation~\ref{eq:bayeslinearnormal} will become 
\begin{equation} 
    f(y|\mathbf{x}, \boldsymbol{\beta}, \sigma^2)\propto \underbrace{(\sigma^2)^{-v/2}\exp (-\frac{vs^2}{2\sigma^2})}_{f(\sigma)}\underbrace{(\sigma^2)^{-\frac{n-v}{2}}\exp \Big (-\frac{1}{2\sigma^2}(\boldsymbol{\beta}-\hat{\boldsymbol{\beta}})^{\top}(\mathbf{x}^{\top}\mathbf{x})(\boldsymbol{\beta}-\hat{\boldsymbol{\beta}})\Big )}_{f(\boldsymbol{\beta}|\sigma)}
\end{equation}
where $vs^2=(\mathbf{y}-\mathbf{x}\hat{\boldsymbol{\beta}})^{\top}(\mathbf{y}-\mathbf{x}\hat{\boldsymbol{\beta}})=SSE$ where $SSE$ is the Error Sum of Squares by definition, such that $v=n_{obs}-p$, $n_{obs}$ is the number of observed observations in the data. 
Then $f(\sigma^2)$ can be written as a proportional density to the inverse gamma distribution,
\begin{equation} 
    f(\sigma^2)\propto (\sigma^2)^{-\frac{v}{2}-1}\exp (-\frac{vs^2}{2\sigma^2})=(\sigma^2)^{-\frac{v}{2}-1}\exp (-\frac{SSE}{2\sigma^2})
\end{equation}
In this $\sigma^2\sim \text{Inverse-Gamma}(\alpha, \phi)$, we have $\alpha = \frac{v}{2}=\frac{n_{obs}-p}{2}$ and $\phi = \frac{1}{2}vs^2=\frac{SSE}{2}$.
From the inverse-gamma property, when 
\begin{equation} 
    \sigma^2\sim \text{Inverse-Gamma}(\frac{n_{obs}-p}{2}, \frac{SSE}{2})
\end{equation}
and assume $\exists\lambda$, such that 
\begin{equation} 
    \sigma^2=\frac{SSE}{2}/\lambda
\end{equation}
then $\lambda$ can be transformed to
\begin{equation} 
    \lambda = \frac{\chi^2_{n_{obs}-p}}{2}
\end{equation}
since $\chi^2_{df}=\text{Gamma}(\frac{df}{2},2)$ where $df$ stands for degrees of freedom, so 
\begin{equation} 
    \sigma^2=\frac{\frac{SSE}{2}}{\frac{\chi^2_{n_{obs}-p}}{2}}=\frac{SSE}{\chi^2_{n_{obs}-p}}
\end{equation}
Thus, we can sample the standard deviation of the missing data distribution from 
\begin{equation} 
    \sigma^*=\hat{\sigma}\sqrt{\frac{SSE}{\chi^2_{n_{obs}-p}}}=\hat{\sigma}\sqrt{\frac{SSE}{g}}
\end{equation}
and from sampling $g\sim \chi^2_{n_{obs}-p}$. 
Moreover, we know 
\begin{equation} 
    \text{Var}(\hat{\boldsymbol{\beta}})=\sigma^2(\mathbf{x}^{\top}\mathbf{x}^{-1})=\mathbf{V}
\end{equation}
so the marginal distribution for $\boldsymbol{\beta}$ is written as 
\begin{equation}\label{eq:betadistmi}
    \boldsymbol{\beta}\sim N(\hat{\boldsymbol{\beta}}, \sigma^2(\mathbf{x}^{\top}\mathbf{x})^{-1})
\end{equation}
Also note that when a random variable $T\sim N(\mathbf{m}, \mathbf{c})$, then $T$ can be generated from a standard normal variable $\mathbf{u}$ with 
\begin{equation} 
    T=\mathbf{m}+\mathbf{L}\mathbf{u}
\end{equation}
where $\mathbf{L}$ is the cholesky decomposition of $\mathbf{c}$ such that $\mathbf{c}=\mathbf{L}\mathbf{L}^{\top}$.
For $\boldsymbol{\beta}$, from~\ref{eq:betadistmi}, it can be derived as 
\begin{align} 
    \boldsymbol{\beta}&=\hat{\boldsymbol{\beta}}+\sigma ( \mathbf{x}^{\top}\mathbf{x} )^{-1/2}\mathbf{u}_1\\
    &=\hat{\boldsymbol{\beta}}+\mathbf{u}_1\mathbf{V}^{-1/2}
\end{align}
where $\mathbf{V}^{-1/2}$ is the cholesky decomposition of $\mathbf{V}$. 
Adjusting for $\sigma^*$ to make sure $\boldsymbol{\beta}$ matches the variability implied by the random draw of $\sigma^*$, 
\begin{equation} 
    \boldsymbol{\beta}^*=\hat{\boldsymbol{\beta}}+\frac{\sigma^*}{\hat{\sigma}}\mathbf{u}_1\mathbf{V}^{-1/2}
\end{equation}
such that $\mathbf{u}_1$ is a row vector of $p$ independent draws from a standard normal distribution, $u_{1k}\stackrel{iid}{\sim} N(0,1)$, and $k=1,\ldots,p$. 
The imputation for $y_i^*$ is then computed as 
\begin{equation} 
    y_i^*=\boldsymbol{\beta}^*\mathbf{x}_i,
\end{equation}
Therefore, the comprehensive steps of the multiple imputation for continuous missing data can be summarized as follows:
\begin{enumerate} 
    \item Calculate $\hat{y}=\hat{\boldsymbol{\beta}}\mathbf{x}$ using $y_{obs}$, and $\hat{\boldsymbol{\beta}}$ can be obtained easily, as well as $\hat{\sigma}$, and $\text{Var}(\hat{\boldsymbol{\beta}})=\mathbf{V}$
    \item Draw $g\sim \chi^2_{n_{obs}-p}$ for one random draw 
    \item Calculate $\sigma^*=\hat{\sigma}/\sqrt{SSE/g}$
    \item Draw a $p$ dimensional vector $\mathbf{u}_1$ such that $u_{1k}\stackrel{iid}{\sim} N(0,1)$ and $k=1,\ldots,p$
    \item Calculate $\boldsymbol{\beta}^*=\hat{\boldsymbol{\beta}}+\frac{\sigma^*}{\hat{\sigma}}\mathbf{u}_1\mathbf{V}^{1/2}$ such that $\mathbf{V}^{1/2}$ is the cholesky decomposition of $\mathbf{V}$
    \item Draw $u_{2i}\stackrel{iid}{\sim} N(0,1)$ 
    \item Impute $y_{mis,i}=\boldsymbol{\beta}^*\mathbf{x}_i$ 
    \item (Option: PMM) Match the imputed value to the nearest observed value.
    \item Repeat 2.\ to 7.\ (2.\ to 8. for PMM) for $M$ times to obtain $M$ complete datasets
\end{enumerate}

The Predictive Mean Matching (PMM) in the above imputation steps is corresponding to one of imputation options proposed by~\citet{rubin1986statistical}, then further discussed by~\citet{little1988missing}. 
This option serves as an algorithm that matches the imputed value to the nearest observed value. 
Mathematically, when we have obtained the $\hat{y}_{mis,i}$ from step 7, we then replace $\hat{y}_{mis,i}$ to $y_{obs,j}$ where $j=\arg\min_{j\in\{1,...,n_{obs}\}} |\hat{y}_{mis,i}-y_{obs,j}|$.
This option arguably preserved the distribution of the observed data because it imputed missing values with actual observed values rather than the predicted ones. 
This helps maintain the variability and distributional property in certain ways of the data. 
But this may suffer from lack of uncertainty and potential biases when the missing proportion is high. 

This imputation method is widely applied to most of missing continuous variable problems. 
However, in genetic epidemiology, when clustered data are in a special form - some continouos variable may be correlated among individuals within one family, this method may yield a biased estimate. 
Because this assumes an independent structure when adding the uncertainty to the imputed value, which ignores the potential correlations that may be due to the genetic nature. 
Therefore, in next section, we modified this imputation approach by adding the kinship correlations to the imputation step, and derived the conditional expectation of the missing value for individual $i$ given other individuals $-i$. 

\section{Multiple Imputation with Kinship Matrix}
\raggedbottom
In order to account for the kinship correlations, the conditional normal distribution needs to be adjusted where the variable contains the missing components are said to be multivariate. 
Denote $\mathbf{y}$ as the continuous variable vector that is subject to missing, $\mathbf{x}$ is other covariates that formed a design matrix that determined to be eligible to be part of the imputation model, and $\boldsymbol{\beta}$ is simply the parameter of the imputation model.
\begin{equation} 
    \mathbf{y}|\mathbf{x},\boldsymbol{\beta}\sim MVN(\boldsymbol{\beta}\mathbf{x}, \sigma_g^2K+\sigma_e^2I)
\end{equation}
where in the linear mixed effect regression form with flexible covariance matrix, the model can be written as 
\begin{equation} 
    \mathbf{y}=\mathbf{x}\boldsymbol{\beta}+\mathbf{u}+\mathbf{e}
\end{equation}
such that two random parts of the model have two multivariate normal distributions,
\begin{equation*}
    \mathbf{u}\sim MVN(0,\sigma_g^2K),~ \text{and}~ \mathbf{e}\sim MVN(0,\sigma_e^2I)
\end{equation*}
where $\sigma_g^2$ introduces the genetic variances, $\sigma_e^2$ introduces the residual variances. 
$K$ is the kinship matrix with the diagonal of 1, and $I$ is the identity matrix. 
Note that the flexible covariance matrix can be written as 
\begin{equation}
    \boldsymbol{\Sigma}=\sigma_g^2K+\sigma_e^2I
\end{equation}
So in this multivariate version of linear mixed effects model, the kinship matrix is being adapted. 
With unknown mean and fully unknown covariance matrix, the prior of the covariance matrix can be selected as an Inverse Wishart distribution with degrees of freedom and the scale parameter in Bayesian sense. 
However, in this case, the covariance matrix is not fully unknown. 
The only unknown parts of this covariance matrix is introduced by simply $\sigma_g^2$ and $\sigma_e^2$. 
Deriving the Bayesian prior of these two parameters is challenging, especially these are incorporated into the partially known multivariate parameter $\boldsymbol{\Sigma}$. 
Therefore, prior definition steps are adjusted to empirical estimates of these two variances. 
Although, in a normal model, the inverse-gamma distribution can be an option as a conjugate prior for variance component. 
The conjugacy simplifies the mathematical derivation, but there are two variance components when the flexible covariance matrix for this designed imputation model, which defers from the regular choice. 
With the above preliminary settings before the imputation step, the adjusted multiple imputation can be concluded into following steps: 
\begin{enumerate} 
    \item Obtain the kinship matrix $K$ among all individuals 
    \item Calculate the estimates of $\hat{y}=\mathbf{x}\boldsymbol{\beta}$, obtain estimates of $\hat{\boldsymbol{\beta}}$, $\hat{\sigma_g}^2$, $\hat{\sigma_e}^2$, $\text{Var}(\hat{\boldsymbol{\beta}})=\mathbf{V}$. In this step, naturally, $\boldsymbol{\hat{\Sigma}}$ is obtained.
    \item Suppose there are $p$ predictors in the imputation model, draw $p$-dimensional vector $\mathbf{w}_1$ such that $w_{1k}\stackrel{iid}{\sim} N(0,1)$ where $k=1,...,p$
    \item Calculate $\boldsymbol{\beta}^*=\hat{\boldsymbol{\beta}}+w_1\mathbf{V}^{1/2}$ such that $\mathbf{V}^{1/2}$ is the cholesky decomposition of $\mathbf{V}$
    \item Obtain $\mu_i^*=\boldsymbol{\beta}^*\mathbf{x}_i$
    \item Obtain the conditional expectations 
    \begin{equation}\label{eq:conditionalexpectation}
        E(y_{mis,i}|\mathbf{y}_{-i})=\mu_i^*+\hat{\boldsymbol{\Sigma}}_{i,-i}\hat{\boldsymbol{\Sigma}}_{-i,-i}^{-1}(\mathbf{y}_{-i}-\boldsymbol{\mu}^*_{-i})
    \end{equation}
    where $\boldsymbol{\mu}_{-i}^*$ is the observed mean vector for other individuals than $i$.
    \item Impute $y_{mis,i}=E(y_{mis,i}|\mathbf{y}_{-i})$
    \item (Option: PMM) Match the imputed value to the nearest observed value
    \item Repeat 3. to 7. (3. to 8. for PMM) for $M$ times to obtain $M$ completed datasets.
\end{enumerate}

The Equation~\ref{eq:conditionalexpectation} holds with the following proof: 
\begin{proof}
    Consider a random vector $\mathbf{x}\sim MVN(\boldsymbol{\mu},\boldsymbol{\Sigma})$, where $\mathbf{x}\in \mathbb{R}^n$, $\boldsymbol{\mu}\in\mathbb{R}^n$ is the mean vector, and $\boldsymbol{\Sigma}\in\mathbb{R}^{n\times n}$ is the covariance matrix.
    We partition $\mathbf{x}$ into $x_i$ and $\mathbf{x}_{-i}$, where $x_i$ is a scalar and $\mathbf{x}_{-i}$ is the $(n-1)$-dimensional vector:
    \begin{equation*}
        \mathbf{x}=\begin{pmatrix} x_i \\ \mathbf{x}_{-i} \end{pmatrix},\\ \boldsymbol{\mu} = \begin{pmatrix} \mu_i \\ \boldsymbol{\mu}_{-i} \end{pmatrix},\quad \boldsymbol{\Sigma}=\begin{pmatrix} \Sigma_{ii} & \boldsymbol{\Sigma}_{i,-i}\\ \boldsymbol{\Sigma}_{-i,i} & \boldsymbol{\Sigma}_{-i,-i} \end{pmatrix}
    \end{equation*}
    Here $\Sigma_{ii}$ is a scalar, which is a variance of $x_i$, $\boldsymbol{\Sigma}_{i,-i}$ is the covariance vector between $x_i$ and $\mathbf{x}_{-i}$, and $\boldsymbol{\Sigma}_{-i,-i}$ is the covariance matrix of $\mathbf{x}_{-i}$. 
    With the property of the multivariate normal distribution, such that the conditional distribution of a subset of a multivariate normal vector given another subset is still normal, with mean and covariance that can be computed from the parameters of the joint distribution. 
    We will only focus on the conditional expectation in this case. 
    The joint density of $(x_i, \mathbf{x}_{-i})$ is given by: 
    \begin{equation*}
        f\left ( \begin{pmatrix} x_i \\ \mathbf{x}_{-i} \end{pmatrix} \right ) = \frac{1}{(2\pi)^{n/2} |\boldsymbol{\Sigma}|^{1/2}} \exp \left ( -\frac{1}{2} \begin{pmatrix} x_i - \mu_i \\ \mathbf{x}_{-i} - \boldsymbol{\mu}_{-i} \end{pmatrix}^T \boldsymbol{\Sigma}^{-1} \begin{pmatrix} x_i - \mu_i \\ \mathbf{x}_{-i} - \boldsymbol{\mu}_{-i} \end{pmatrix} \right )
    \end{equation*}
    Now we compute the inverse of the covariance matrix $\boldsymbol{\Sigma}$, using the block matrix inversion formula: 
    \begin{equation*}
        \boldsymbol{\Sigma}^{-1} = \begin{pmatrix} \Sigma_{ii} & \boldsymbol{\Sigma}_{i,-i} \\ \boldsymbol{\Sigma}_{-i,i} & \boldsymbol{\Sigma}_{-i,-i} \end{pmatrix}^{-1} = \begin{pmatrix} A & B \\ C & D \end{pmatrix}
    \end{equation*}
    where the submatrices $A$, $B$, $C$, and $D$ are given by:
    \[
        A = (\Sigma_{ii} - \boldsymbol{\Sigma}_{i,-i} \boldsymbol{\Sigma}_{-i,-i}^{-1} \boldsymbol{\Sigma}_{-i,i})^{-1}
    \]
    \[
        B = -A \boldsymbol{\Sigma}_{i,-i} \boldsymbol{\Sigma}_{-i,-i}^{-1}
    \]
    \[
        C = -\boldsymbol{\Sigma}_{-i,-i}^{-1} \boldsymbol{\Sigma}_{-i,i} A
    \]
    \[
        D = \boldsymbol{\Sigma}_{-i,-i}^{-1} + \boldsymbol{\Sigma}_{-i,-i}^{-1} \boldsymbol{\Sigma}_{-i,i} A \boldsymbol{\Sigma}_{i,-i} \boldsymbol{\Sigma}_{-i,-i}^{-1}
    \]  
    The quadratic form in the exponent of the joint density is
    \begin{equation}\label{eq:quadraticform}
        \begin{pmatrix} x_i - \mu_i \\ \mathbf{x}_{-i} - \boldsymbol{\mu}_{-i} \end{pmatrix}^T \boldsymbol{\Sigma}^{-1} \begin{pmatrix} x_i - \mu_i \\ \mathbf{x}_{-i} - \boldsymbol{\mu}_{-i} \end{pmatrix}
    \end{equation}
    Expanding~\ref{eq:quadraticform}, and using the symmetric of the covariance matrix that we know $C=B^{\top}$, we can get 
    \begin{align*}
        &\begin{pmatrix} x_i - \mu_i \\ \mathbf{x}_{-i} - \boldsymbol{\mu}_{-i} \end{pmatrix}^T \boldsymbol{\Sigma}^{-1} \begin{pmatrix} x_i - \mu_i \\ \mathbf{x}_{-i} - \boldsymbol{\mu}_{-i} \end{pmatrix} \\
        &= (x_i - \mu_i)^T A (x_i - \mu_i) + (x_i - \mu_i)^T B (\mathbf{x}_{-i} - \boldsymbol{\mu}_{-i}) + \\
        &+(\mathbf{x}_{-i} - \boldsymbol{\mu}_{-i})^T C (x_i - \mu_i) + (\mathbf{x}_{-i} - \boldsymbol{\mu}_{-i})^T D (\mathbf{x}_{-i} - \boldsymbol{\mu}_{-i}) \\
        &= (x_i - \mu_i)^T A (x_i - \mu_i) + 2 (x_i - \mu_i)^T B (\mathbf{x}_{-i} - \boldsymbol{\mu}_{-i}) + (\mathbf{x}_{-i} - \boldsymbol{\mu}_{-i})^T D (\mathbf{x}_{-i} - \boldsymbol{\mu}_{-i})
    \end{align*}
    The conditional mean comes directly from the term \( 2 (x_i - \mu_i)^T B (\mathbf{x}_{-i} - \boldsymbol{\mu}_{-i}) \), leading to the linear adjustment term \(\boldsymbol{\Sigma}_{i,-i} \boldsymbol{\Sigma}_{-i,-i}^{-1} (\mathbf{x}_{-i} - \boldsymbol{\mu}_{-i})\), where
    \begin{equation*}
        E(x_i|\mathbf{x}_{-i})=\mu_i+\boldsymbol{\Sigma}_{i,-i}\boldsymbol{\Sigma}_{-i,-i}^{-1}(\mathbf{x}_{-i}-\boldsymbol{\mu}_{-i})
    \end{equation*}
\end{proof}
Using the conditional expectation captures the genetic correlations when imputing the missing variable. 
In fact, once the correlation matrix is known, even broader scope than the genetic epidemiology, can still apply this proposed multiple imputation method.

Generally, when there is no preliminary knowledge on the dataset, one should use as many variables as possible for the imputation model covariate~\cite{rubin2018multiple}.
Also, in the survival analysis, there has been always an argument on whether the $\log(T)$ or $\log(H_0(T))$ should be used within the imputation model~\cite{white2009imputing}. 
In this thesis, we have used both $\log(T)$ and $\log(H_0(T))$ in the simulation study to determine the differences to avoid potential biases caused by misspecification of the model. 

\section{Pooling Step and Variance Estimation}
After multiple imputation, Rubin's rules are applied to combine the results from multiple datasets.
These rules derive overall estimates and associated variances that incorporate the uncertainty due to missing data.
Specifically, Rubin's rule for variance estimation involves calculating within-imputation and between-imputation variances.
The total variance is the sum of the average within-imputation variance and the between-imputation variance, adjusted by a factor related to the number of imputations.
This approach ensures that the variability due to missing data is properly reflected in the final statistical inferences.
Rubin's rule is advantageous compared to empirical variance estimation because it accounts for both within-imputation and between-imputation variances, providing a more accurate reflection of the uncertainty due to missing data.
This dual consideration enhances the validity of statistical inferences by incorporating the variability introduced through the imputation process. 
In contrast, empirical variance estimation may underestimate this uncertainty, leading to overly optimistic conclusions~\cite{rubin2018multiple}.

Based on Rubin's rule, the pooling step of the analysis is defined as
\begin{equation} 
    \bar{\theta}=\frac{1}{M}\sum_{i=1}^M\theta_i
\end{equation}
where $\theta_i$ is the parameter estimate for the $i$-th imputation after $M$ imputations. 
Thus, $\bar{\theta}$ represents the pooled parameter estimates.
According to \citet{barnard1999miscellanea}, variance estimation is structured at two levels: within-imputation variance and between-imputation variance.
The within imputation variance is similar to the pooled parameter estimation, where 
\begin{equation} 
    V_W=\frac{1}{M}\sum_{i=1}^MSE_i^2
\end{equation}
such that $V_W$ determines the within-imputation variance, and it's simply ``pooled" variance among all analyses from the imputed dataset $i$. 
The between-imputation variance $V_B$, on the other hand, accounts for the extra variances caused by the missing data is defined as 
\begin{equation} 
    V_B=\frac{\sum_{i=1}^M(\theta_i-\bar{\theta})^2}{M-1}
\end{equation}
where $V_B$ is the unbiased estimate of between-imputation variance. 
Given that $\theta$ is estimated using the only the finite $M$ imputed datasets,~\citet{van2018flexible} note that $V_B$ is the approximation of the variance when $M\rightarrow\infty$.
Therefore, the total variance can be written as
\begin{equation} 
    V_{\text{Total}}=V_W+V_B+\frac{V_B}{M}
\end{equation}
when $M$ is large enough, $V_{\text{Total}}$ tends to primarily consist of within-imputation and between-imputation variances.
The test statistic (Wald statistic) for the pooled estimate is 
\begin{equation}\label{eq:teststatFdist}
    \frac{(\bar{\theta}-\theta_0)^2}{V_{\text{Total}}}\sim F_{1,df_{adj}}
\end{equation}
where $\theta_0$ is the null hypothesis for the estimate, which~\ref{eq:teststatFdist} can be derived to
\begin{equation} 
    \frac{\bar{\theta}-\theta_0}{\sqrt{V_{\text{Total}}}}\sim t_{df_{adj}, \frac{1-\alpha}{2}}
\end{equation}
such that $\alpha$ is the significance level. 
Define 
\begin{equation}
    \rho=\frac{V_B+\frac{V_B}{M}}{V_{\text{Total}}}
\end{equation}
and 
\begin{equation}
    r=\frac{V_B+\frac{V_B}{M}}{V_W}
\end{equation}
where $\rho$ is the fraction of missing information, which quantifies the proportion of the total variance that is attributable to the fact that the data have been imputed $M$ times. 
Also, $r$ is the relative increase in variance due to non-responses (missing data), which measures the increase in variance due to the missing data compared to if there were no missing data. 
In the interpretation, $\rho$ gives the proportion of the total variance that is due to the uncertainty introduced by the missing data. 
A higher $\rho$ indicates a larger fraction of the total uncertainty comes from the fact that the data were imputed.
A higher $r$ indicates that the missing data have a larger impact on the overall variance. 
In the first MI literature by~\citet{rubin1987multiple}, the degrees of freedom was defined as
\begin{equation}\label{eq:dfold}
    df_{old}=(M-1)\times (1+\frac{1}{r})^2
\end{equation}
The estimated degrees of freedom for the observed data, adjusted for the missing information, are
\begin{equation}\label{eq:dfobs}
    df_{obs}=\frac{(n-k)+1}{(n-k)+3}(n-k)(1-\rho)
\end{equation}
where $n$ is the number of observations in each imputed data, and $k$ is the number of parameters to be estimated.
\citet{barnard1999miscellanea} further alternated the calculation of the degrees of freedom by combining Equations~\ref{eq:dfold} and~\ref{eq:dfobs}
\begin{equation}
    df_{adj}=\frac{df_{old}df_{obs}}{df_{old}+df_{obs}}
\end{equation}
When conducting the statistical tests on the pooled estimates, $df_{adj}$ is used for the t-distribution degrees of freedom. 
The confidence interval is straightforward, where we obtain 
\begin{equation}
    SE_{pooled}=\sqrt{V_{\text{Total}}}
\end{equation}
then we simply calculate 
\begin{equation}
    CI=\bar{\theta}\pm t_{df_{adj},\frac{1-\alpha}{2}}\times SE_{pooled}
\end{equation}

In the multiple imputations, Rubin's rules provide a robust framework for combining estimates from multiple imputed datasets, effectively accounting for both within-imputation and between-imputation variances. 
This approach enhances the validity of statistical inferences by appropriately reflecting the uncertainty due to missing data, while the structure of statistical testing remains straightforward without complex alternations.
%% *** Include chapter files here. ***


%% This adds a line for the Bibliography in the Table of Contents.
\addcontentsline{toc}{chapter}{Bibliography}
%% *** Set the bibliography style. ***
%% (change according to your preference/requirements)
\bibliographystyle{unsrtnat}
%% *** Set the bibliography file. ***
%% ("thesis.bib" by default; change as needed)
\bibliography{westernthesis.bib}

%% *** NOTE ***
%% If you don't use bibliography files, comment out the previous line
%% and use \begin{thebibliography}...\end{thebibliography}.  (In that
%% case, you should probably put the bibliography in a separate file and
%% `\include' or `\input' it here).

\end{document}
